% Part: normal-modal-logic
% Chapter: syntax-and-semantics
% Section: language-modal-logic

\documentclass[../../../include/open-logic-section]{subfiles}

\begin{document}

\olfileid{nml}{syn}{lan}

\olsection{मूलभूत मोडल तर्कशास्त्राची भाषा}

\begin{defn}
मोडल तर्कशास्त्राच्या मूलभूत भाषेत पुढील गोष्टी असतात:
\begin{enumerate}
  \tagitem{prvFalse}{!!{falsity} साठीचा विधानीय स्थिरांक~$\lfalse$.}{}
  \tagitem{prvTrue}{!!{truth} साठीचा विधानीय स्थिरांक~$\ltrue$.}{}
  \item !!{propositional variable}s चा !!{denumerable}s संच: $\Obj
    p_0$, $\Obj p_1$, $\Obj p_2$, \dots
  \item विधानीय संयोजक: \startycommalist
  \iftag{prvNot}{\ycomma $\lnot$ (निषेधन)}{}%
  \iftag{prvAnd}{\ycomma $\land$ (संधी)}{}%
  \iftag{prvOr}{\ycomma $\lor$ (विकल्प)}{}%
  \iftag{prvIf}{\ycomma $\lif$ (!!{conditional})}{}%
  \iftag{prvIff}{\ycomma $\liff$ (!!{biconditional})}{}.
  \tagitem{prvBox}{मोडल परिकर्मी $\Box$.}{}
  \tagitem{prvDiamond}{मोडल परिकर्मी $\Diamond$.}{}
\end{enumerate}
\end{defn}

\begin{defn}
मूलभूत मोडल भाषेतील \emph{!!^{formula}s} पुढीलप्रमाणे
  विगमनाने परिभाषित होतात:
\begin{enumerate}
\tagitem{prvFalse}{$\lfalse$ हे आण्विक !!{formula} आहे.}{}

\tagitem{prvTrue}{$\ltrue$ हे आण्विक !!{formula} आहे.}{}

\item प्रत्येक !!{propositional variable} $\Obj p_i$ हे (आण्विक) !!{formula} आहे.

\tagitem{prvNot}{$!A$ हे !!a{formula} असेल, तर $\lnot !A$
  हे !!a{formula} आहे.}{}

\tagitem{prvAnd}{$!A$ आणि $!B$ ही !!{formula}s असतील, तर $(!A \land
  !B)$ हे !!a{formula} आहे.}{}

\tagitem{prvOr}{$!A$ आणि $!B$ ही !!{formula}s असतील, तर $(!A \lor !B)$
  हे !!a{formula} आहे.}{}

\tagitem{prvIf}{$!A$ आणि $!B$ ही !!{formula}s असतील, तर $(!A \lif !B)$
  हे !!a{formula} आहे.}{}

\tagitem{prvIff}{$!A$ आणि $!B$ ही !!{formula}s असतील, तर $(!A \liff !B)$
  हे !!a{formula} आहे.}{}

\tagitem{prvBox}{$!A$ हे !!a{formula} असेल, तर $\Box !A$
  हे !!a{formula} आहे.}{}

\tagitem{prvDiamond}{$!A$ हे !!a{formula} असेल, तर $\Diamond !A$
  हे !!a{formula} आहे.}{}

\tagitem{limitClause}{यांशिवाय दुसरे काहीही !!a{formula} नाही.}{}
\end{enumerate}
\end{defn}

\begin{tagblock}{defNot,defOr,defAnd,defIf,defIff,defTrue,defFalse,defBox,defDiamond}
\begin{defn}
परिभाषित परिकर्मी वापरून बनवलेली सूत्रे पुढीलप्रमाणे
समजायची:

\begin{tagenumerate}{defTrue,defFalse,defNot,defOr,defAnd,defIf,defIff,defBox,defDiamond}

\tagitem{defTrue}{$\ltrue$ हे पुढीलचे संक्षिप्त रूप आहे:
  \iftag{prvFalse}{$\lnot\lfalse$}{$(!A \lor \lnot !A)$;
    येथे एखादे निश्चित आण्विक !!{formula}~$!A$ घेतले आहे}.}{}

\tagitem{defFalse}{$\lfalse$ हे पुढीलचे संक्षिप्त रूप आहे:
  \iftag{prvTrue}{$\lnot\ltrue$}{$(!A \land \lnot !A)$;
    येथे एखादे निश्चित आण्विक !!{formula}~$!A$ घेतले आहे}.}{}

\tagitem{defNot}{$\lnot !A$ हे $!A \lif \lfalse$ चे संक्षिप्त रूप आहे.}{}

\tagitem{defOr}{$!A \lor !B$ हे पुढीलचे संक्षिप्त रूप आहे:
  \iftag{prvAnd}{$\lnot(\lnot !A \land \lnot !B)$}{$\lnot !A \lif
    !B$}.}{}

\tagitem{defAnd}{$!A \land !B$ हे पुढीलचे संक्षिप्त रूप आहे:
  \iftag{prvOr}{$\lnot(\lnot !A \lor \lnot !B)$}{$\lnot (!A \lif
    \lnot !B)$}.}{}

\tagitem{defIf}{$!A \lif !B$ हे पुढीलचे संक्षिप्त रूप आहे:
  \iftag{prvOr}{$(\lnot !A \lor !B)$}{$\lnot (!A \land \lnot !B)$}.}{}

\tagitem{defIff}{$!A \liff !B$ हे $(!A \lif !B) \land (!B
  \lif !A)$ चे संक्षिप्त रूप आहे.}{}

\tagitem{defBox}{$\Box !A$ हे $\lnot\Diamond\lnot !A$ चे संक्षिप्त रूप आहे. }{}

\tagitem{defDiamond}{$\Diamond !A$ हे $\lnot\Box\lnot !A$ चे संक्षिप्त रूप आहे.}{}
\end{tagenumerate}
\end{defn}
\end{tagblock}

!!a{formula}~$!A$ मध्ये $\Box$ किंवा $\Diamond$ नसेल,
तर त्याला \emph{मोडलतारहित} म्हणतो.

\end{document}
